% Shared current index algorithm. All added displays are unnumbered.
\clearpage
\sessionheading{Disconnected gauge sectors and the factorization identity}
\label{sec:disconnected}
A product gauge group is not necessarily a product of independent theories:
matter charged under several factors couples those factors. For each nonzero
hypermultiplet $h$, define its gauge support
\[
 S_h=\{a:\ R_{h,a}\ne\mathbf1\}.
\]
\code{_split_disconnected_sectors} builds an undirected NetworkX graph with
one vertex per gauge factor. For a support $(a_1,\ldots,a_r)$, it inserts
the $r-1$ edges $(a_1,a_2),\ldots,(a_1,a_r)$. A star joins every vertex of
that support by a path; a complete clique would give the same connected
components. Overlapping supports join transitively. Zero-multiplicity matter
is ignored and isolated gauge vertices are retained. Connected components are
obtained with \code{networkx.connected_components} \cite{networkx-components}.

For components $C_1,\ldots,C_s$, each charged hyper belongs to exactly one
component. The code preserves factor order, restricts representation and beta
contribution dictionaries, and converts a one-factor product hyper into
\code{HyperData}. Multi-factor components retain \code{ProductHyperData}.
No discarded spectator singlet contributes an extra dimension factor.

Let $L_\alpha$ denote the letters of component $\alpha$ and $L_0$ all matter
with empty gauge support. The exponent splits because Adams substitution and
summation are linear, so the plethystic identity gives
\[
 L=L_0+\sum_\alpha L_\alpha,\qquad
 \PE[L]=\PE[L_0]\prod_\alpha\PE[L_\alpha].
\]
Normalized Haar integration also factorizes, since each $L_\alpha$ depends
only on $g_\alpha\in G_{C_\alpha}$:
\[
 \begin{aligned}
 \mathcal I
 &=\int\prod_\alpha d\mu_\alpha(g_\alpha)\,
        \PE[L_0]\prod_\alpha\PE[L_\alpha(g_\alpha)]\\
 &=\mathcal I_0\prod_\alpha\mathcal I_\alpha,\qquad
 \mathcal I_0=\PE[L_0],\quad
 \mathcal I_\alpha=\int d\mu_\alpha\,\PE[L_\alpha].
 \end{aligned}
\]
This follows directly from the plethystic definition and product measure
\cite{plethystics,gpr}; the graph is the implementation's test for these
independent integration variables.

\textbf{Why all-gauge-singlet matter appears once.} A gauge-singlet hyper
has empty support. Dropping it loses a physical free-matter factor. Copying
it into each of $s$ components would replace $\mathcal I_0$ by
$\mathcal I_0^s$. The code puts all such hypers in one additional sector
with no gauge factors, preserving each original multiplicity. ``Once'' means
one assignment of each input species, not changing its multiplicity to one.

\clearpage
\sessionheading{Exact multiplication, flavor characters and database reuse}
At finite cutoff, multiplication and truncation commute in the needed sense:
\[
 \left[\prod_{\alpha=0}^{s}\mathcal I_\alpha\right]_N
 =\left[\prod_{\alpha=0}^{s}[\mathcal I_\alpha]_N\right]_N.
\]
All sector indices have nonnegative $t$ degree. Discarded powers above $N$
cannot multiply later terms back down to $N$. Laurent powers of $y,u$ do not
alter that argument. \code{_truncate_index} filters the exact Sage monomial
dictionary inclusively after each sector and each product. An identical sector
within the same call is looked up or calculated once, but its index is still
multiplied once per occurrence.

\textbf{Flavor refinement.} The same identity holds with external flavor
variables $x$: $\mathcal I(t,y,u;x)=\mathcal I_0(t,y,u;x)
\prod_\alpha\mathcal I_\alpha(t,y,u;x)$. Shared flavor variables are coefficients,
not integrated gauge variables. Products of their characters describe tensor
products and need not be decomposed first. For example, a free full hyper with
$U(1)$ fugacity $x$ contributes $\PE[f_H(x+x^{-1})]$ once. Gauging a shared
flavor subgroup introduces a common Haar integral and can connect sectors;
one cannot then multiply separately projected indices. The current code is
flavor-unrefined: this is the extension's mathematical contract, not an
implemented flavor API.

\textbf{Optional theory database.} \code{calculate_index} and
\code{calculate_index_internal} accept \code{theory_db_connection=None}.
For each gauged sector, \code{common.n2_theory_db.find_superconformal_index}
matches the existing canonical Lagrangian hash and reads an index only if it
is a JSON string with recorded \code{superconformal_index_order >= order}.
Canonicalization retains the established full/half and conjugation conventions;
factor IDs are ignored but input factor order remains significant.
Missing/null indices and unknown/lower cutoffs are misses. A malformed index
or one with negative $t$ powers is rejected by the caller and recalculated.
The cutoff is never inferred from the highest surviving monomial.

Misses use \code{_calculate_sector_index}: character basis, degree-bounded FORM
expansion/cache and the existing LiE singlet projection. Free sectors bypass
the theory lookup. No sector theories or indices are inserted. A supplied
connection is borrowed read-only and is never committed, rolled back or closed
by this calculation; it must belong to the calling process/thread and already
have the current schema. GUI/database workers pass their own connection.
Without one, sector splitting and multiplication still apply.

The character-cache file is a distinct parameter:
\code{char_cache_database_path=None}. It replaces \code{cache_directory} and
\code{database_path} in the index API; the index and database-worker CLI use
\code{--char-cache-database}. The independent FORM file uses
\code{form_cache_database_path}. The standalone character-cache class/builder
retain their own directory/file API and CLI. Saved GUI key
\code{cache/character_database} is unchanged.

\clearpage
\sessionheading{Product-group bottlenecks and recorded evidence}
The 14 September investigation found 39 full-index FORM timeouts in a
32-worker run through $t^{18}$: 29 inputs had disconnected gauge sectors.
The main isolated cost was constructing the joint formal-character exponential,
before Python parsing, LiE projection or database writes. A printing-disabled
control for disconnected $SU(3)\times SU(3)$ theory 9780 still took 516.16 s
and retained 2,649,141 formal terms, whereas the final projected index has
168 monomials. This locates that delay in symbolic computation, not an
observed database deadlock. Large expansions also incur parsing and memory
costs. The report's saved timeout setting was 60 s; the code default was 600 s.
Those are dated observations, not assertions about the user's current settings.

\begin{tabularx}{\textwidth}{@{}>{\raggedright\arraybackslash}p{39mm}Y@{}}
\toprule Implementation & Effect and evidence\\\midrule
Disconnected-sector split & Projects small independent sectors before
multiplying their final polynomials. On 15 September, theory 9780 took
13.753 s with one unique sector calculation and exact agreement with all
168 saved combined-result monomials. Character cache was a private populated
copy; FORM cache was empty; no MySQL connection was used.\\
Degree-bounded FORM & Helps connected sectors too, by omitting products
above the remaining degree budget. The 16 September matched bifundamental
run is recorded in the FORM explanation.\\
Remaining limits & Connected sectors can still have many formal monomials
and expensive LiE intermediates. Concurrent misses can still duplicate one
FORM program; no in-progress job claim or streaming projection was added.\\\bottomrule
\end{tabularx}

\textbf{Evidence files.} \code{output/benchmarks/product_form_20260914/report.md}
retains the original profiling controls and inputs;
\code{output/benchmarks/disconnected_sectors_20260915.json} records the
implemented split. The earlier diagnostic degree-bounded variant became the
production generator on 16 September. Historical timings with different
printing or cache conditions are not interchangeable benchmarks.

\textbf{Regression coverage.} \code{test/test_disconnected_index.py} tests
transitive and trifundamental connectivity, isolated vectors, singlet matter,
zero multiplicities, factor order, repeated sectors and exact comparison with
an unsplit calculation. Database cases cover missing/null indices, unknown,
lower, equal and higher cutoffs, canonical identity and borrowed-connection
ownership. The recorded 15 September combined suite passed 353 tests in
58.615 s with no skips. The subsequent cache-argument revision passed 355
tests in 60.129 s. These are implementation-time results, not reruns performed
for this document update.
